\documentclass[mestrado,onehalfspacing,impressao]{UnB-PPGT}

% Definir metadados obrigatórios
\titulo{Teste do Sistema de Páginas Pré-textuais}
\autor{João}{Silva}
\orientador{Prof. Dr. Maria Santos}{Universidade de Brasília}
\datadefesa{15}{12}{2024}
\coordenador{Prof. Dr. Carlos Oliveira}{Universidade de Brasília}

% Definir palavras-chave e keywords
\palavraschave{transporte, mobilidade, sustentabilidade, LaTeX}
\keywords{transport, mobility, sustainability, LaTeX}

% Definir membros da banca
\membrobanca{Prof. Dr. Ana Costa}{Universidade de Brasília}
\membrobanca{Prof. Dr. Pedro Lima}{Universidade Federal do Rio de Janeiro}

\begin{document}

% Teste dos ambientes pré-textuais

\begin{dedicatoria}
À minha família, pelo apoio incondicional durante toda esta jornada acadêmica.
\end{dedicatoria}

\begin{agradecimentos}
Gostaria de expressar minha sincera gratidão ao meu orientador, Prof. Dr. Maria Santos, pela orientação dedicada e pelos valiosos conselhos ao longo desta pesquisa.

Agradeço também aos colegas do Programa de Pós-Graduação em Transportes pelo companheirismo e pelas discussões enriquecedoras que contribuíram para o desenvolvimento deste trabalho.

Por fim, agradeço à Universidade de Brasília pela oportunidade de realizar esta pesquisa e pelo ambiente acadêmico estimulante.
\end{agradecimentos}

\begin{resumo}
Este trabalho apresenta um estudo sobre o desenvolvimento de um modelo LaTeX para dissertações e teses do Programa de Pós-Graduação em Transportes da Universidade de Brasília. O objetivo principal é criar um template que siga fielmente as normas institucionais e facilite a elaboração de documentos acadêmicos de alta qualidade.

A metodologia adotada baseou-se na análise do template Word oficial do PPGT e na adaptação das melhores práticas dos modelos LaTeX existentes na UnB. O resultado é um sistema integrado que automatiza a geração das páginas pré-textuais, formatação de elementos textuais e integração com sistemas de referências bibliográficas.

Os resultados demonstram que o modelo desenvolvido atende aos requisitos estabelecidos, proporcionando uma ferramenta eficiente para a elaboração de trabalhos acadêmicos no âmbito do PPGT-UnB.
\end{resumo}

\begin{abstract}
This work presents a study on the development of a LaTeX template for dissertations and theses of the Graduate Program in Transportation at the University of Brasília. The main objective is to create a template that faithfully follows institutional standards and facilitates the preparation of high-quality academic documents.

The methodology adopted was based on the analysis of the official Word template of PPGT and the adaptation of best practices from existing LaTeX models at UnB. The result is an integrated system that automates the generation of pre-textual pages, formatting of textual elements and integration with bibliographic reference systems.

The results demonstrate that the developed model meets the established requirements, providing an efficient tool for the preparation of academic works within the scope of PPGT-UnB.
\end{abstract}

\chapter{Introdução}

Este é um capítulo de teste para verificar se o sistema de páginas pré-textuais está funcionando corretamente.

\end{document}